\chapter{Arquitectura de microservicios}
\section{Catálogo de APIs verificables}Las rutas siguientes provienen de controladores y routers; los prefijos FastAPI se obtienen de sus routers.
\begin{longtable}{p{.13\textwidth}p{.12\textwidth}p{.37\textwidth}p{.3\textwidth}}\caption{Catálogo resumido de APIs}\toprule Servicio&Método&Ruta&Propósito\\\midrule\endfirsthead\toprule Servicio&Método&Ruta&Propósito\\\midrule\endhead Catalog&GET&/api/catalog/home, /search&Home y búsqueda.\\Catalog&GET/POST&/api/catalog/movies, /\{publicId\}, /session-info&Catálogo y detalle para sesión.\\Identity&POST&/register, /login&Alta y autenticación JWT.\\Community&CRUD&/api/v1/users, /clubs, /watch-rooms&Perfiles, clubes y salas.\\Interaction&GET/POST/PUT/DELETE&/api/v1/reviews, /movies/:id/like, /users/:id/watchlist&Interacciones documentales.\\Cinema&POST/PATCH/GET&/api/v1/sessions, /:id/playback, /:id/state, /:id/chat&Sesión activa.\\Analytics&GET&/top-watched, /kpis, /movies/best-rated&Consultas Athena.\\\bottomrule\end{longtable}
\section{Catalog Service}
Catalog concentra el catálogo cinematográfico. Su código Java Spring Boot se organiza en controladores \texttt{CatalogController}, \texttt{MovieController} y \texttt{HealthController}, servicios, repositorios JPA y entidades. Una petición sigue Controller $\rightarrow$ Service $\rightarrow$ Repository $\rightarrow$ PostgreSQL; los DTO separan la respuesta pública de entidades. \texttt{Movie} es la entidad principal y se relaciona con géneros, artistas, fuentes de video y subtítulos; \texttt{MovieVideoSource} y \texttt{MovieSubtitle} pertenecen al detalle reproducible. \texttt{MovieRepository} usa \texttt{EntityGraph} para cargar relaciones requeridas. La configuración JDBC apunta a PostgreSQL y el contenedor expone 8081. En IaC, Catalog recibe API Gateway $\rightarrow$ VPC Link $\rightarrow$ ALB interno 8081 $\rightarrow$ TGCatalog $\rightarrow$ ambas EC2 Prod. Un GET de película atraviesa ese recorrido, busca por UUID público mediante controlador/servicio/repositorio y devuelve JSON. No se verificó un middleware de autorización propio en los controladores auditados.
\section{Identity Service}
Identity es la autoridad de cuentas y tokens. FastAPI registra \texttt{/register} y \texttt{/login}; \texttt{models.py} define \texttt{users} con UUID, email único, nombre, hash y rol. En registro, Pydantic valida el cuerpo, el servicio consulta la unicidad, aplica hash de contraseña y persiste SQLAlchemy. En login, consulta usuario, verifica hash y \texttt{core/security.py} emite JWT HS256 con \texttt{sub}, \texttt{exp} e \texttt{iss=identity-service}. El secreto y la expiración proceden de configuración, no se reproducen valores. El contenedor escucha 9000; IaC define IdentityAPI, listener 9000 y TGIdentity con las dos EC2.\section{Community Service}
Community persiste la capa social: \texttt{users} locales asociados al \texttt{sub} de Identity, \texttt{clubs}, \texttt{memberships}, \texttt{watch\_rooms} y \texttt{watch\_participants}. Un club tiene propietario, visibilidad y miembros; Membership vincula usuario-club con rol OWNER/ADMIN/MEMBER; WatchRoom identifica película, anfitrión y club opcional; Participant registra nickname y rol HOST/VIEWER. \texttt{deps.py} valida Bearer JWT, exige issuer y puede crear el perfil local en el primer acceso autenticado. Por ello Community no almacena contraseñas ni emite tokens. Rutas verificables incluyen GET/PUT usuarios, CRUD clubes, altas/listas de membresía y POST/GET/join de watch rooms. Usa FastAPI, SQLAlchemy y MySQL en 8000.
\section{Interaction Service}
Interaction implementa escritura social de alta variabilidad con Node.js, Express, Mongoose y MongoDB. \texttt{auth.js} extrae Bearer, verifica JWT con algoritmo configurado y coloca identidad para rutas protegidas. Review contiene \texttt{movie\_id}, \texttt{user\_id}, score 0--10, comentario, fecha y snapshots movie/user; Like contiene identificadores, fecha y snapshots. Watchlist e historial usan el mismo patrón de contexto embebido. Los índices por movie, user y el compuesto único movie/user en Review y Like impiden duplicados. El recorrido de un like es HTTP $\rightarrow$ authenticate $\rightarrow$ movie route $\rightarrow$ controlador $\rightarrow$ modelo Mongoose $\rightarrow$ MongoDB; reseña, lista e historial siguen la misma separación. Las rutas incluyen crear/listar/editar/borrar reseñas, like/unlike, métricas, watchlist e historial; los endpoints de ingesta están protegidos por clave específica.
\section{Cinema Session Service}
La rama auditada es \texttt{feature/http-polling-chat}. Community persiste una watch room; Cinema administra el estado activo transitorio de una sesión sincronizada. \texttt{session\_service.py} crea sesiones, consulta película en Catalog y consulta room en Community; mantiene datos en memoria. Las rutas permiten crear/listar/obtener/borrar sesiones, añadir/quitar/listar participantes, PATCH playback, GET state y POST chat. El cliente puede obtener estado por HTTP polling mediante GET state; además \texttt{ws.py} define WebSocket para eventos y chat. Ambos mecanismos coexisten: no se describe como exclusivamente polling. Cinema consulta Interaction para estadísticas/likes con timeout de dos segundos en las llamadas visibles. Un PATCH playback actualiza el estado local; otro cliente puede recibir evento WebSocket o consultar GET state. Con dos EC2 detrás de ALB, solicitudes consecutivas pueden caer en réplicas distintas: si no comparten memoria, una sesión creada en Prod1 puede no existir en Prod2.\section{Data Analytics Service}
Analytics no accede al OLTP en sus rutas: FastAPI delega a \texttt{app/athena/client.py}. Este crea cliente boto3, envía \texttt{StartQueryExecution}, sondea el estado y devuelve filas cuando alcanza SUCCEEDED; FAILED y CANCELLED se tratan como terminales. \texttt{ATHENA\_OUTPUT\_BUCKET} configura resultados. Rutas verificables cubren top watched, actor/género, mejor calificación, likes, actividad social, KPIs, embudo, retención y salud de clubs. El contenedor expone 8002 y la plantilla define AnalyticsAPI/TGAnalytics.
